\documentclass[11pt]{article}
\usepackage{varnernotes}

\newcommand{\E}{\mathbb{E}}
\newcommand{\Cov}{\operatorname{Cov}}
\newcommand{\R}{\mathbb{R}}

\begin{document}

\lecturenotes{5a}{Multivariate GBM and Dirichlet Portfolio Weights}{Fall 2026}{September 22, 2026}{Jeffrey D. Varner}

\learningobjectives
  {Model correlated assets}
  {Write and solve a multiple asset GBM model whose covariance factor generates a specified log-return covariance rate.}
  {Estimate covariance}
  {Construct, annualize, and validate a sample covariance matrix while distinguishing covariance from correlation.}
  {Sample portfolio weights}
  {Use a Dirichlet distribution to explore the long-only allocation simplex without confusing random sampling with optimization.}

\section{Why Separate Asset Models Are Insufficient}

Single-asset GBM models can be simulated independently, but an independently simulated portfolio erases co-movement. Firms respond together to market news, sectors share inputs, and hedges move in opposing directions. Portfolio uncertainty depends not only on each asset's volatility but also on every pairwise covariance.

Let $M\geq2$ be the number of assets. Collect the prices in $\mathbf{S}_t\in\R_{>0}^M$, arithmetic GBM drifts in $\boldsymbol\mu\in\R^M$, and independent Wiener processes in $\mathbf{W}_t\in\R^M$. A multiple asset GBM (MAGBM) uses a loading matrix (covariance factor) $A\in\R^{M\times M}$ (units years$^{-1/2}$):
\begin{equation}
  \label{eq:mgbm-sde}
  \frac{dS_{i,t}}{S_{i,t}}
  =\mu_i\,dt+\sum_{\ell=1}^{M}A_{i\ell}\,dW_{\ell,t},
  \qquad i\in\{1,\ldots,M\}.
\end{equation}
The diffusion covariance rate is $C:=AA^{\mathsf T}$, whose entries have units years$^{-1}$. Its diagonal entries satisfy $C_{ii}=\sigma_i^2$, while $C_{ij}$ records instantaneous log-return covariance per unit time. Applying It\^o's lemma to $\log S_{i,t}$ gives the expected continuously compounded growth rate
\begin{equation}
  \label{eq:drift-growth-relation}
  \bar g_i:=\E[g_{i,t,\Delta t}]
  =\mu_i-\frac12 C_{ii},
  \qquad
  \mu_i=\bar g_i+\frac12 C_{ii}.
\end{equation}
Thus $\mu_i$ is reserved for the arithmetic drift in the price SDE, while $g_{i,t,\Delta t}:=(\Delta t)^{-1}\log(S_{i,t+\Delta t}/S_{i,t})$ is the realized growth rate. Both have units years$^{-1}$, but they are different quantities.

\begin{figure}[ht]
  \centering
  \begin{tikzpicture}[x=1.35cm,y=0.86cm,>=Latex,font=\sffamily\small]
    \node[draw=vnrule,fill=white,rounded corners=1pt,
      minimum width=2.35cm,minimum height=0.78cm] (z) at (0,0) {independent $\mathbf Z$};
    \node[draw=vnink,fill=vnpanel,rounded corners=1pt,
      minimum width=2.25cm,minimum height=0.78cm] (a) at (3.25,0) {covariance factor $A$};
    \node[draw=vnrule,fill=white,rounded corners=1pt,
      minimum width=2.65cm,minimum height=0.78cm] (az) at (6.5,0) {correlated shock $A\mathbf Z$};
    \node[draw=vnink,fill=vnpanel,rounded corners=1pt,
      minimum width=2.55cm,minimum height=0.78cm] (s) at (10.0,0) {joint price update};
    \foreach \u/\v in {z/a,a/az,az/s}
      \draw[vnlink,line width=1.0pt,-{Latex[length=3mm]}] (\u)--(\v);
    \node[below,text=vnmuted] at (3.25,-0.62) {$AA^{\mathsf T}=C$};
    \node[below,text=vnmuted] at (6.5,-0.62) {$\Cov(A\mathbf Z)=C$};
  \end{tikzpicture}
  \caption{A covariance factor $A$ with $AA^{\mathsf T}=C$ converts independent standard-normal shocks into shocks with covariance $C$. The same shock vector must drive all assets on one simulated path to preserve their joint dependence.}
  \label{fig:correlated-shocks}
\end{figure}

\section{Exact Transition and Covariance}

The logarithm of each component receives its own It\^o correction, while the off-diagonal covariance enters through the common shock vector (\propref{mgbm-transition}).

\begin{proposition}{Exact constant-parameter MAGBM transition}{mgbm-transition}
Let $A\in\R^{M\times M}$, $C:=AA^{\mathsf T}$, $\boldsymbol\mu\in\R^M$, and $\Delta t>0$ years. Define $\bar g_i:=\mu_i-C_{ii}/2$. For independent standard normal shocks $\mathbf Z\sim\mathcal N(\mathbf0,I_M)$, drawn fresh at every step and shared by all assets within the step, the exact one-step solution of \eqnref{mgbm-sde} is
\begin{equation}
  \label{eq:mgbm-transition}
  S_{i,t+\Delta t}=S_{i,t}\exp\!\left[
  \bar g_i\Delta t
  +\sqrt{\Delta t}\,(A\mathbf Z)_i
  \right].
\end{equation}
Define the dimensionless log-return vector $\mathbf r_t:=\log(\mathbf S_{t+\Delta t}\oslash\mathbf S_t)$ and the realized growth-rate vector $\mathbf g_t:=\mathbf r_t/\Delta t$. Then
\begin{equation}
  \label{eq:mgbm-log-covariance}
  \mathbf g_t=\bar{\mathbf g}+\frac{1}{\sqrt{\Delta t}}A\mathbf Z,
  \qquad
  \E[\mathbf g_t]=\bar{\mathbf g}=\boldsymbol\mu-\frac12\operatorname{diag}(C),
  \qquad
  \Cov(\mathbf r_t)=C\Delta t,
  \qquad
  \Cov(\mathbf g_t)=\frac{C}{\Delta t}.
\end{equation}
\end{proposition}

Equation \eqnref{mgbm-log-covariance} gives the scaling rule. Writing $\Sigma_r:=\Cov(\mathbf r_t)$ and $\Sigma_g:=\Cov(\mathbf g_t)$,
\begin{equation}
  \label{eq:three-covariances}
  \Sigma_r=(\Delta t)^2\Sigma_g,
  \qquad
  C=\frac{\Sigma_r}{\Delta t}=\Delta t\,\Sigma_g.
\end{equation}
The three matrices are not interchangeable: $\Sigma_g$ is appropriate for a growth-rate portfolio objective, while $C$ drives GBM shocks. Multiplying a daily log-return covariance by ``252'' estimates $C$; it does not estimate $\Sigma_g$.

\section{Estimating and Validating the Covariance Matrix}

Let $G\in\R^{N\times M}$ store $N\geq2$ observed growth-rate vectors by rows (time periods $k=1,\ldots,N$, aligned and equally spaced), and let $\mathbf g'\in\R^M$ be the vector of sample means (the sample mean $g'$ of L3a and L4b, one per asset; $\bar{\mathbf g}$ remains the population mean $\boldsymbol\mu-\tfrac12\operatorname{diag}(C)$). With $\mathbf1_N$ the length-$N$ vector of ones, define the centered matrix $\widetilde G:=G-\mathbf1_N\mathbf g'^{\mathsf T}$. The unbiased growth-rate covariance is
\begin{equation}
  \label{eq:sample-covariance}
  \widehat\Sigma_g:=
  \frac{1}{N-1}\widetilde G^{\mathsf T}\widetilde G.
\end{equation}
It is symmetric and positive semidefinite: for any $\mathbf x\in\R^M$, $\mathbf x^{\mathsf T}\widehat\Sigma_g\mathbf x=\lVert\widetilde G\mathbf x\rVert_2^2/(N-1)\geq0$, and $\widehat C=\Delta t\,\widehat\Sigma_g$ inherits both properties. For GBM simulation, use $\widehat C$; its diagonal reproduces the L4b volatilities, $\widehat\sigma_i=\sqrt{\widehat C_{ii}}=\sigma_{g,i}\sqrt{\Delta t}$.

Covariance depends on units and volatility scale. Correlation normalizes those effects:
\begin{equation}
  \label{eq:correlation}
  \rho_{ij}:=\frac{\widehat\Sigma_{g,ij}}
  {\sqrt{\widehat\Sigma_{g,ii}\widehat\Sigma_{g,jj}}},
  \qquad -1\leq\rho_{ij}\leq1
\end{equation}
when both variances are positive; $\widehat\Sigma_g$, $\widehat\Sigma_r$, and $\widehat C$ share the same correlations because they differ by positive scalars. Zero correlation means zero linear covariance under the chosen sample and horizon; it does not generally imply statistical independence.

A Cholesky factor is convenient when $\widehat\Sigma_g$ is positive definite; the symmetric square root $C^{1/2}$ is another valid covariance factor, and any $A$ with $AA^{\mathsf T}=C$ works. Duplicated assets, exact linear dependence, or too few observations make it only semidefinite (the centered matrix has rank at most $N-1$, so $M\geq N$ guarantees singularity), causing an unpivoted Cholesky routine to fail. Eigenvalue or singular-value square roots can handle semidefinite matrices. Arbitrarily adding diagonal ``jitter'' changes the model and must be documented rather than hidden.

Sampling error is often the dominant problem. Pairwise estimates can shift across windows, and a matrix with $M(M+1)/2$ distinct entries becomes difficult to estimate when $M$ is large relative to $N$. Shrinkage and factor models trade some flexibility for greater stability.

\notebooklink
  {L5a: Computing and Analyzing a Covariance Matrix}
  {https://github.com/varnerlab/CHEME-5660-CourseRepository-Fall-2026/blob/main/lectures/week-5/L5a/CHEME-5660-L5a-Example-CovarianceMatrix-Fall-2026.ipynb}
  {Center a multi-asset growth-rate matrix, calculate $\widehat\Sigma_g$, convert it to the GBM covariance rate $\widehat C$, and verify numerical invariants.}

\section{Buy-and-Hold Portfolio Wealth}

Let total initial wealth be $W_0>0$ dollars and let $\mathbf w\in\R^M$ contain initial dollar fractions. For a fully invested long-only portfolio, $w_i\geq0$ and $\mathbf1^{\mathsf T}\mathbf w=1$. Fractional-share holdings are $n_i=w_iW_0/S_{i,0}$. If no rebalancing occurs, portfolio wealth is
\begin{equation}
  \label{eq:buy-hold-wealth}
  W_t=\sum_{i=1}^{M}n_iS_{i,t}
  =W_0\sum_{i=1}^{M}w_i\frac{S_{i,t}}{S_{i,0}}.
\end{equation}
The initial weights generally drift as relative prices change: the fraction of wealth in asset $i$ at time $t$ is $w_i(t)=n_iS_{i,t}/W_t$. A constant-weight portfolio is a different, dynamically rebalanced strategy with additional turnover and cost assumptions.

\notebooklink
  {L5b: Correlated MAGBM Portfolio Simulation}
  {https://github.com/varnerlab/CHEME-5660-CourseRepository-Fall-2026/blob/main/lectures/week-5/L5b/CHEME-5660-L5b-Example-MAGBM-Portfolio-Fall-2026.ipynb}
  {Generate correlated asset paths, construct a buy-and-hold portfolio, and compare simulated wealth with observed and benchmark wealth.}

\section{Dirichlet Sampling Explores the Simplex}

The feasible long-only weight set is the $(M-1)$-dimensional simplex $\Delta^{M-1}:=\{\mathbf w\in\R^M:w_i\geq0,\ \sum_iw_i=1\}$. The Dirichlet distribution generates random points on its interior and is a useful design-of-experiments tool (\defnref{dirichlet}).

\begin{definition}{Dirichlet portfolio weights}{dirichlet}
Let $\boldsymbol\alpha=(\alpha_1,\ldots,\alpha_M)$ with every concentration parameter $\alpha_i>0$, and let $\alpha_0:=\sum_{i=1}^{M}\alpha_i$. A random vector $\mathbf W\sim\operatorname{Dirichlet}(\boldsymbol\alpha)$ lies in the interior of $\Delta^{M-1}$ and has component mean, variance, and covariance
\begin{equation}
  \label{eq:dirichlet-moments}
  \E[W_i]=\frac{\alpha_i}{\alpha_0},
  \qquad
  \operatorname{Var}(W_i)=
  \frac{\alpha_i(\alpha_0-\alpha_i)}
  {\alpha_0^2(\alpha_0+1)},
  \qquad
  \operatorname{Cov}(W_i,W_j)=-\frac{\alpha_i\alpha_j}{\alpha_0^2(\alpha_0+1)}\ (i\neq j).
\end{equation}
For every $\alpha_i=1$, the distribution is uniform with respect to volume on the simplex; it is not uniform in each component separately (each marginal is $\operatorname{Beta}(1,M-1)$).
\end{definition}

Small equal concentrations place more mass near sparse corners, while large equal concentrations cluster weights near equal allocation. Unequal concentrations encode a directional sampling preference. None of these draws is ``optimal'' unless an objective and constraints are subsequently evaluated. Dirichlet sampling explores candidate portfolios; L5b solves a minimum-variance optimization problem.

\notebooklink
  {L5a: Sampling Portfolio Weights with the Dirichlet Distribution}
  {https://github.com/varnerlab/CHEME-5660-CourseRepository-Fall-2026/blob/main/lectures/week-5/L5a/CHEME-5660-L5a-Example-Dirichlet-PortfolioWeights-Fall-2026.ipynb}
  {Draw weight vectors on the three-asset simplex for several concentration vectors, map the growth-rate mean and variance of random five-asset portfolios, and track buy-and-hold wealth and weight drift.}

\begin{figure}[ht]
  \centering
  \begin{tikzpicture}[x=1cm,y=1cm,font=\sffamily\small]
    \coordinate (a) at (0,0); \coordinate (b) at (7,0); \coordinate (c) at (3.5,4.7);
    \draw[vnink,line width=0.9pt] (a)--(b)--(c)--cycle;
    \node[below left] at (a) {$w_1=1$};
    \node[below right] at (b) {$w_2=1$};
    \node[above] at (c) {$w_3=1$};
    \foreach \x/\y in {1.0/.45,2.1/.7,2.8/1.4,3.5/.75,4.2/1.45,5.15/.55,3.5/2.25,2.65/2.25,4.35/2.15,3.5/3.45}
      \fill[vnlink] (\x,\y) circle (1.8pt);
    \node[text=vnmuted,align=center] at (3.5,-0.75) {$w_i\geq0$ and $w_1+w_2+w_3=1$};
  \end{tikzpicture}
  \caption{The three-asset long-only simplex. Dirichlet draws provide valid candidate allocations throughout the triangle; their density depends on the concentration vector $\boldsymbol\alpha$.}
  \label{fig:simplex}
\end{figure}

\newpage
\section*{Optional Advanced Material}

\notebooklink
  {L5a Advanced: Sampling Error and Shrinkage in Covariance Estimation}
  {https://github.com/varnerlab/CHEME-5660-CourseRepository-Fall-2026/blob/main/lectures/week-5/L5a/advanced/covariance-estimation/CHEME-5660-L5a-Advanced-CovarianceEstimation-Fall-2026.ipynb}
  {Compare the eigenvalue spectrum of the sample covariance with the Marchenko--Pastur law, shrink the estimate toward a structured target, and measure the effect on a minimum-variance portfolio out of sample.}

\notebooklink
  {L5a Advanced: Rolling Correlations}
  {https://github.com/varnerlab/CHEME-5660-CourseRepository-Fall-2026/blob/main/lectures/week-5/L5a/advanced/rolling-correlation/CHEME-5660-L5a-Advanced-RollingCorrelation-Fall-2026.ipynb}
  {Track pairwise and average correlations through 2014 to 2024 with rolling and exponentially weighted windows, and quantify what a constant covariance rate misses in stressed markets.}

\newpage
\thispagestyle{fancy}
\keytakeaways
  {In this lecture, we coupled GBM assets through a covariance matrix, built correlated shocks from independent ones with a covariance factor, estimated the covariance from growth rates, and used Dirichlet draws to explore valid portfolio weights.}
  {Correlated assets need a shared shock and a covariance factor}
  {One shock vector per step, mixed by $A$ with $AA^{\mathsf T}=C$, gives log returns with covariance $C\Delta t$; each asset keeps its own correction $\bar g_i=\mu_i-C_{ii}/2$.}
  {Covariance is estimated from centered growth rates and scaled by the step}
  {$\widehat\Sigma_g$ is symmetric and positive semidefinite by construction, $\widehat C=\Delta t\,\widehat\Sigma_g$ drives the simulation, and singularity, sampling error, and window dependence must be checked before use.}
  {Sampling is not optimization}
  {Dirichlet draws explore the long-only simplex; concentration parameters control where the draws land but do not encode a minimum-risk objective.}
  {Next, turn estimated means and covariance into the minimum-variance portfolio and the efficient frontier, and simulate the resulting portfolios with correlated MAGBM paths.}

\begin{readinglist}
  \item Harry Markowitz, \href{https://www.jstor.org/stable/2975974}{``Portfolio Selection,''} \emph{Journal of Finance} 7(1), 77--91 (1952).
  \item Robb J. Muirhead, \emph{Aspects of Multivariate Statistical Theory} (1982), for covariance matrices and multivariate normal transformations.
  \item Thomas M. Cover and Joy A. Thomas, \emph{Elements of Information Theory}, 2nd ed. (2006), Appendix, for the Dirichlet family and simplex-valued distributions.
\end{readinglist}

\vfill
{\sffamily\footnotesize\color{vnmuted}\textbf{Educational use and risk notice.} These notes are for instruction only and do not constitute investment advice. Financial decisions involve risk, and model outputs depend on assumptions.}

\end{document}
